\chapter{Introdução}
\label{cap:introducao}

O poquêr, particularmente a modalidade \textit{Texas Hold'Em}, é considerado um dos maiores desafios atuais para a inteligência artificial, pois envolve aspectos vitais como a tomada de decisão diante da incerteza, gerenciamento de risco e a necessidade de entender o comportamento dos oponentes. Ao contrário de jogos tradicionais como xadrez e \textit{Go}, onde todas as informações estão disponíveis para os jogadores, o poquêr apresenta a característica de informação incompleta, exigindo decisões baseadas em dados parciais e sob a influência de elementos aleatórios. Essa particularidade tem incentivado o desenvolvimento de novas técnicas em \textit{machine learning}, com foco especial em aprendizado de máquina e modelagem probabilística.

Neste cenário, os métodos que utilizam redes neurais artificiais têm se tornado cada vez mais significativos, sendo capazes de identificar padrões complexos a partir de extensas quantidades de dados, sem depender diretamente do conhecimento anterior do programador. O processo de treinamento desses modelos, que pode ser supervisionado, não supervisionado ou por reforço, possibilita a identificação de táticas de jogo, a previsão de jogadas dos adversários e a realização de escolhas otimizadas, fundamentais na mecânica do poquêr.

Contudo, a complexidade do espaço de estados no \textit{Texas Hold'Em}, devido ao enorme número de combinações possíveis de cartas, apostas e estilos de jogo dos oponentes, torna impraticável a exploração de todas as opções disponíveis através de métodos convencionais. Para contornar esse desafio, destacam-se os algoritmos de \textit{Monte Carlo}, incluindo o \gls{MCTS} e o \gls{MCCFR}, que conseguem analisar amostras selecionadas das possibilidades do jogo, estimando chances de ganhar, retorno esperado e ajustando estratégias em tempo real. Essa metodologia simula diversos cenários possíveis, utilizando redes neurais para determinar os melhores caminhos, mesmo com limitações computacionais.

As mais recentes inovações neste campo envolvem a combinação de métodos de \textit{Monte Carlo} com arquiteturas de redes neurais profundas, como as redes \textit{actor-critic}, permitindo que o agente se desenvolva através de autojogo e levando a políticas que se aproximam do equilíbrio de Nash, um padrão de referência em jogos de soma zero. Exemplos como \textit{DeepStack}, \textit{Libratus} e \textit{Pluribus} evidenciam a eficácia dessas ferramentas ao vencer campeões mundiais em competições reais, revelando táticas sólidas para partidas entre dois jogadores e também em ambientes com múltiplos competidores.

Dessa forma, a criação de uma rede neural capaz de jogar poquêr \textit{Texas Hold’Em} utilizando técnicas de \textit{Monte Carlo} fornece uma contribuição significativa para o entendimento da IA aplicada a decisões complexas. Este estudo tem como objetivo executar modelagem teórica até práticas em ambientes simulados, demonstrando como a união entre aprendizado profundo e simulações probabilísticas pode resultar em agentes de alto desempenho em cenários com informações incompletas e vários adversários.
